\subsection[Challenge 23: One-loop correction of quasi-PDF]{Challenge 23: One-loop correction of quasi-PDF}
\textbf{Problem ID:} \texttt{Challenge\_23\_main}\\
\textbf{Primary inferred discipline:} High Energy \& Nuclear Physics\\
\textbf{Python implementation:} \texttt{challenges/23/answer.py}

\subsubsection*{Problem}
\paragraph*{Problem setup}
In perturbative quantum chromodynamics (QCD), one always encounters divergent Feynman diagrams in loop corrections. In the Large Momentum Effective Theory (LaMET) framework for calculating parton physics, the perturbative matching coefficient is calculated from Feynman diagrams that involves one external quark or gluon with momentum $p$, which is set to be on-shell and massless, i.e., $p^2=0$ and
\begin{align}
    I &= \int \frac{d^d k}{(2\pi)^d} \frac{N(k,p)}{D(k,p)} \,,
\end{align}
where $d$ is the dimension of the Minkowski space time with the metric $g^{\mu\nu}=\mbox{diag}\{1,-1,\ldots,-1\}$.

There are two types of divergences, ultraviolet (UV) and infrared (IR). Before one carries out the loop integration over $k^\nu=(k^0, k^1, \ldots, k^{d-1})$, it is necessary to first identify the UV and IR divergences of each diagram through the superficial degree of divergence (SDD).

\emph{UV divergences.} In the limit of $k^\nu \to \infty$ for each $\nu=0,1,2,\ldots,d-1$, the SDD is $d-t$ if $N(k,p)/D(k,p)\rightarrow 1/k^t$, where $t$ is an integer.

\emph{IR divergences.} There are two types of IR divergences, soft and collinear divergences.
\begin{itemize}[leftmargin=*]
\item Soft divergence: In the limit of $k^\nu \to 0$ for each $\nu=0,1,2,\ldots,d-1$, $N(k,p)/D(k,p)\rightarrow 1/k^t$;
\item Collinear divergence: In the limit of $k^\nu \to \lambda p^\nu$, with $0<\lambda< 1$,  $N(k,p)/D(k,p)\rightarrow 1/(k - \lambda p)^t$.
\end{itemize}

One way to regulate these divergent integrals is dimensional regularization, where the space-time dimension is set to be $d=4-2\epsilon$ with a continuous parameter $\epsilon$. It works under the assumption that the regulated results can be analytically continued in $\epsilon$. The UV divergences are regulated with $\epsilon_{\rm UV}>0$, and the IR divergences are regulated with $\epsilon_{\rm IR}<0$.

\paragraph*{Main problem}

The so-called sail-diagram contribution to the one-loop matrix element of the quasi-PDF in a positive-energy quark state with momentum $p^\nu=(p^z,0,0,p^z)$ is
\begin{align}
    \tilde q_{\rm sail}(x,p^z,\epsilon,\mu) &=(\mu^2)^\epsilon \frac{\exp(\epsilon\gamma_E)}{(4\pi)^\epsilon}\int_{-\infty}^\infty \frac{dz }{2\pi}e^{ixp^z z} \int \frac{d^d k}{(2\pi)^d} \frac{k^0+k^z}{k^2 (p-k)^2 (p^z-k^z)}e^{-ik^z z}.
\end{align}
Carry out the above integral in three intervals: $x < 0$, $0<x<1$, and $x>1$. Then expand the final result in $\epsilon$ to $O(\epsilon^0)$ in the $\overline{\rm MS}$ scheme. The final results of $\tilde q_{\rm sail}(x,p^z,\epsilon,\mu)$ should be expressed as functions in $x$, $p^z$, $\epsilon_{\rm UV}$ or $\epsilon_{\rm IR}$, and $\mu$.
